% HD-G31. GPT-6 Astra Pro /root/g34. 21 November 2026.
% Inherited effort; numerical setting not exposed.
\subsection{HD-G31: the primitive normal rank-three branch}
\label{sec:rank-three-primitive-scope}

Retain the actual finite coefficient group, its prescribed natural
primitive Jacobian factor, and the imaginary-quadratic field and
polarization conventions of HD-S41--HD-S43.  In particular let
$G\subset GL_m(K)$ be finite and absolutely irreducible with exact
character field $K$, and let its full sphere-cover factor be a Weil
sixfold with period rank nine.  Consider the remaining case in which
no nonidentity element of $G$ has moving rank one or two.  Let $N_3$
be the subgroup generated by all actual elements $x$ with
$\operatorname{rank}_K(x-1)=3$.

The branch where $N_3$ is absolutely irreducible, complex primitive,
and tensor indecomposable has
\[
                             15\le m\le35.                 \tag{G31.1}
\]
Here tensor indecomposable means that its representation does not
preserve a fixed nontrivial tensor product; the published theorem
below also treats permutations of tensor factors.  In this branch,
for $K=\mathbb Q(\sqrt{-d})$ with
$d\in\{2,7,11,19,43,67,163\}$, every actual moving-rank-three
element is an involution.  For $K=\mathbb Q(i)$ its nonidentity
spectrum is one of
\[
                  (-1,-1,-1),\qquad(i,i,-i),\qquad(-i,-i,i).
                                                               \tag{G31.2}
\]
Combining (G31.1) with HD-R12 excludes this branch altogether for
$d=163$, and restricts it to $21\le m\le35$ for $d=43$ and
$33\le m\le35$ for $d=67$.  No realization or exclusion of the
remaining finite ranges is asserted here.

\paragraph{1. Exact published threshold and all exceptional degrees.}
The input is Guralnick--Saxl, \emph{Generation of finite almost
simple groups by conjugates}, J.~Algebra \textbf{268} (2003),
519--571, Theorem~7.1 on pp.~548--549 and its shorter
Theorem~7.A on p.~521.  For a finite irreducible primitive
and tensor-indecomposable complex linear group they define
\[
 \nu_G(V)=\min_{g,\lambda:\,\lambda g\ne1}
                         \operatorname{rank}(\lambda g-1).
\]
Outside their stated alternatives the exact bound is
$\nu_G(V)>\max\{2,\sqrt m/2\}$.  The alternatives are
finite classical groups in their defining-characteristic natural
modules; the scalar deleted-permutation alternative; the listed
small defining-characteristic modules of degree at most eight;
the cross-characteristic exceptions of degree at most ten;
and extraspecial or tensor-cube exceptions of degree eight.
In characteristic zero and degree at least36 only the scalar
deleted-permutation alternative remains.  In that alternative
the permutation degree is $s=m+1$, since in characteristic zero
the constant line splits off and the deleted module has dimension
$s-1$.  Theorem~7.A explicitly retains multiplication by scalars.

The entire original theorem and exception list were read from the
publicly indexed original article text.  The radical's placement
is also explicitly corroborated in the published primary paper
J.~I.~Hall, \emph{Periodic simple groups of finitary linear
transformations}, Ann.~of Math.~\textbf{163} (2006), 445--498,
printed p.~480, immediately before Theorem~5.2; p.~481 records
$k(d)=4d^2$ for $d\ge3$.  Both original Hall pages were visually
read.  The original GS PDF image was not retrieved, and no visual
read of that PDF is claimed.  The inclusive endpoint $m=36$ uses
the strict inequality in the original GS theorem, rather than
strengthening Hall's separately stated $m>k(d)$ formulation.

The hypotheses apply to the actual $G$.  A $G$-block system would
also be an $N_3$-block system, and irreducibility makes $N_3$
transitive on it.  Similarly a fixed $G$-tensor product would be
one for $N_3$.  Thus $G$ is primitive and tensor indecomposable.
There is an actual rank-three element, since $N_3$ is irreducible
in degree greater than one.  It gives $\nu_G(V)\le3$.
For $m\ge36$ the published theorem therefore applies with
$3\le\sqrt m/2$.

\paragraph{2. A literal rank-two commutator in the original group.}
The scalar deleted-module conclusion supplies an actual complex
intertwiner $C_0$ and the rational standard augmentation matrices
$P_\pi$ on
\[
     V_0=\{(z_1,\ldots,z_s):\textstyle\sum z_i=0\},
      \qquad
     C_0\rho(g)C_0^{-1}=a_gP_{\pi(g)},\qquad
                   \pi(G)=A_s\ \hbox{or}\ S_s.             \tag{G31.3}
\]
The pair $(a_g,\pi(g))$ is unique.  A permutation scalar on
$V_0$ preserves every line $\mathbb C(e_i-e_j)$ and hence every
unordered pair of labels; for $s\ge3$ it fixes all labels and
its scalar is one.  Thus both factors in (G31.3) are actual
homomorphisms.

Take the two even permutations
\[
                 u=(12)(45),\qquad v=(23)(45).
\]
Choose actual lifts $g,h\in G$ under $\pi$.  Scalar factors
cancel in the actual commutator, giving
\[
 C_0\rho([g,h])C_0^{-1}
      =[P_u,P_v]=P_{(132)}.                                \tag{G31.4}
\]
The $s$-coordinate permutation $(132)$ fixes $s-2$ dimensions;
removing the invariant constant line leaves fixed dimension
$s-3$ in the $(s-1)$-dimensional augmentation module.  Its moving
rank is exactly two.  Since $\rho([g,h])-1$ is a matrix over
$K$, its rank remains two before extension to $\mathbb C$.
This is an actual element of the original group, contradicting
the residual rank condition.  No character equality is being
used as a substitute for this commutator map.

Finally, for the actual sphere-cover local system, the parabolic
dimension formula is
\[
                  6=-2m+\sum_{j=1}^{b}c_j.
\]
The rank-nine period assumption requires $b-3\ge9$; here each
$c_j\ge3$.  Consequently $2m+6\ge36$ and $m\ge15$.
Together with (G31.4) this proves (G31.1).

\paragraph{3. The exact field, normalizer, and original corner bridge.}
For completeness the model in the contradiction has an actual
$K$-intertwiner and the same primitive-factor maps; it need not
be treated merely as a complex-character description.  Commutators
in (G31.3) are scalar free, and their permutation images generate
$A_s$.  Hence
\[
                 C_0\rho(G')C_0^{-1}=P(A_s).                \tag{G31.5}
\]
Fix this literal isomorphism $\psi:G'\to A_s$, use the rational
basis $e_j-e_s$, and solve the finite linear equations
\[
                  C\rho(x)=P_{\psi(x)}C\quad(x\in G')      \tag{G31.6}
\]
over $K$.  The standard $A_s$ augmentation module is absolutely
irreducible: its permutation action is two-transitive, so the
permutation-module commutant has dimension two; the constant
line splits and leaves an irreducible augmentation module.
The complex solution space of (G31.6) is therefore one dimensional
and contains $C_0$.  Gaussian elimination has the same rank over
$K$ and $\mathbb C$, so a nonzero $K$-solution exists and is a
nonzero scalar multiple of $C_0$.  In particular it is invertible.

Here is the normalizer map, including arbitrary original
overgroups.  On $V_0$ put $Q=\sum z_j^2$ and $F=\sum z_j^3$.
Their invariant lines are the complete spaces of invariant
quadratic and cubic forms for $A_s$: a monomial of degree at
most three leaves two unused labels, on which a transposition
corrects the parity of a permutation; the rational projection
to $V_0$ then gives the asserted restriction of the symmetric
polynomial spaces.  Every normalizer preserves both lines.
The Hessian of $F$ at $z$ is $6\sum z_jw_jt_j$.  If $z$ has
$k<s$ nonzero entries its rank on $V_0$ is exactly $k$; if all
entries are nonzero its rank is $s-1$ or $s-2$, from the equations
$z_jw_j=\text{constant}$.  Thus the nonzero rank-two Hessian locus
is exactly the union of the lines $\mathbb C(e_i-e_j)$.
Orthogonality for $Q$ says that two such lines have disjoint
label pairs.  The maximal pairwise-intersecting families of at
least four pairs are precisely the full stars of one label.
The normalizer consequently induces a permutation of the labels.
After removing that permutation, a linear map fixes every root
line.  Writing its action as $e_i-e_s\mapsto a_i(e_i-e_s)$
and using the line of $(e_i-e_s)-(e_j-e_s)$ gives $a_i=a_j$.
It is scalar.  This proves the literal normalizer identity
\[
                  N_{GL(V_0)}(A_s)=\mathbb C^*P(S_s).
                                                               \tag{G31.7}
\]
Since $G'$ is characteristic, every $C\rho(g)C^{-1}$ is in
(G31.7).  Multiplying it by the rational inverse permutation
matrix shows that its scalar lies in $K$.  Finiteness makes
that scalar a root of unity.  Thus the original group has the
actual faithful correlated map
\[
                         G\longrightarrow\mu(K)\times S_s.
                                                               \tag{G31.8}
\]
For $\mu(K)=\{\pm1\}$ all its traces are rational, also
contradicting the exact original character field.  In the
Gaussian case (G31.8) is exactly the actual hypothesis of
HD-G26, with no independence of its two projections assumed.

Let $h$ be the original positive $K$-Hermitian coefficient form
and $h_0$ the rational standard form.  For a unique
$a\in\mathbb Q_{>0}$,
\[
                       h_0(Cx,Cy)=a h(x,y).                 \tag{G31.9}
\]
Indeed the two invariant Hermitian forms are proportional by
absolute irreducibility, and their $K$-valued matrices make
the real proportionality factor rational.  If $w$ is the
specified original primitive vector, retaining $Cw$ makes
$h(w,\rho(g)w)/h(w,w)$ literally unchanged.  Thus its normalized
rational matrix-coefficient idempotent in the same relabelled
$\mathbb Q[G]$ is unchanged, with the identical Jacobian image,
restricted theta class, and corner $K$-action.  At the coefficient
comparison level the Schur factors are
$c=2|G|/(m h(w,w))$ and $c_0=c/a$, so the actual Jacobian form
has the same multiplier $c_0a=c$.  The coefficient multiplier
$a$ alone is not being asserted as a Jacobian polarization
multiplier.

For a subsequent change of primitive vector $w$ to $v$ put
$q_w=h(w,w)$, $q_v=h(v,v)$,
\[
 T_{v,w}(x)=v h(w,x)/q_w,\qquad
 b_{v,w}=\frac{2m}{|G|}\sum_{g\in G}
       \operatorname{Re}\operatorname{tr}
          (T_{v,w}\rho(g^{-1}))\,g.                        \tag{G31.10}
\]
Actual matrix-coefficient averaging gives
$b_{v,w}=e_vb_{v,w}e_w$ and
$b_{v,w}^{\dagger}b_{v,w}=(q_v/q_w)e_w$, acting as zero
on the other rational blocks.  The character and its conjugate
are distinct because the original exact field is imaginary
quadratic.  For an integer $M$ clearing the displayed
coefficients, the covariant Jacobian action of $Mb_{v,w}$ is
an actual $K$-linear isogeny; its pullback polarization multiplier
is $M^2q_v/q_w$, with rational inverse given by the adjoint and
the same scalar.  This includes the change to
$v=se_1-\boldsymbol1$ in HD-G26.  That theorem's actual
point-stabilizer and scalar-kernel quotient curves, norm/pullback
maps, and self-adjoint complementary projector then apply to
this very factor.  Every positive rational polarization scalar
$t$ changes the six-dimensional Hermitian determinant by
$t^6=N_{K/\mathbb Q}(t^3)$ and preserves period rank through the
fixed actual isogeny.  These maps explain the complete geometric
bridge even though (G31.4) already excludes the present branch
when $m\ge36$.

\paragraph{4. An actual junior-element exclusion in all remaining degrees.}
A second published input is Guralnick--Tiep, \emph{A problem of
Koll\'ar and Larsen on finite linear groups and crepant
resolutions}, J.~Eur.~Math.~Soc.~\textbf{14} (2012), 605--657,
Theorem~1.4, p.~607, with its proof in Section~5.6,
p.~652.  It states that a finite irreducible complex linear
group of degree at least11 generated, up to scalars, by elements
of age at most one contains an actual complex bireflection of
order two or three.  Its conclusion is an actual element in
the group, not only a projective scalar class.

Here age is defined for the actual matrix: write its eigenvalues as
$\exp(2\pi i a_j)$ with $0\le a_j<1$ and set
$\operatorname{age}(g)=\sum_j a_j$.  No minimization over
external scalars is used.  The projective quantity often denoted
$\operatorname{age}^*$ is not an input to this argument.

There is no noncentral element of age at most one in our actual
$G$.  To prove this, let $J$ be the normal subgroup generated by
all such elements and suppose one exists.  Complete reducibility
and normality give an isotypic decomposition for $J$ permuted by
$G$.  Primitivity forces a single isotypic type, so
$V_{\mathbb C}=W\otimes\mathbb C^e$ with $J$ acting only on
$W$.  Its matrix algebra is $\operatorname{End}(W)\otimes I_e$.
The original $G$ normalizes this algebra.  Choosing the images
of a system of matrix units shows explicitly that each such
normalizer has the form $A_g\otimes B_g$: after removing
$A_g$, it centralizes the full first matrix algebra and hence
acts on the second factor.  If $1<\dim W,e$, this is a fixed
nontrivial tensor product for $G$, already excluded.  If
$\dim W=1$, then $J$ is scalar, contrary to its noncentral
generator.  Hence $e=1$ and $J$ is irreducible.  It is generated
by the specified actual junior elements, so the cited theorem
produces an actual bireflection in $J\subset G$, a contradiction.
All normal-subgroup and tensor maps here are on the actual
complexification; no descent of a separate character is needed
to retain the resulting element of the original $G$.

\paragraph{5. Exact rank-three spectra over the specified fields.}
Let $x\in G$ have moving rank three.  All three nonidentity
eigenvalues have the same order $\ell$.  Otherwise for the
order of one of them, a power of $x$ would fix a nonempty proper
subset of the three moving eigenlines and would have moving rank
one or two.  The moving subspace itself is the actual
$K$-subspace $\operatorname{im}_K(x-1)$, so its characteristic
polynomial is a degree-three polynomial over $K$ whose roots
are primitive $\ell$th roots of unity.

An irreducible factor of $\Phi_\ell$ over $K$ has degree
$\varphi(\ell)$, or degree $\varphi(\ell)/2$ when
$K\subset\mathbb Q(\zeta_\ell)$.  All these factor degrees
are equal, because the cyclotomic extension is Galois abelian.
The common degree must divide three.  Degree one gives exactly
the roots of unity already in $K$; degree three requires
$\varphi(\ell)=6$ and $K$ to be the quadratic subfield of
$\mathbb Q(\zeta_\ell)$.  The latter orders are
$\ell=7,9,14,18$, with quadratic fields respectively
$\mathbb Q(\sqrt{-7})$ or $\mathbb Q(\sqrt{-3})$.
For the field embeddings, put
$\tau=\zeta_7+\zeta_7^2+\zeta_7^4$.  Direct multiplication
using $1+\zeta_7+\cdots+\zeta_7^6=0$ gives
$\tau+\bar\tau=-1$ and $\tau\bar\tau=2$, so
$(2\tau+1)^2=-7$.  In $\mathbb Q(\zeta_9)$ put
$\omega=\zeta_9^3$; then $(2\omega+1)^2=-3$.
These give the literal embeddings of the two quadratic fields.
The groups $(\mathbb Z/7)^*$ and $(\mathbb Z/9)^*$ are
cyclic of order six, so each has a unique subgroup of index two.
The orders14 and18 give the same respective cyclotomic fields.
Equivalently, the actual Galois stabilizer of a primitive root
has index three over $K$, while its index over $\mathbb Q$
is six.

For the seven non-Gaussian, non-Eisenstein fields in the statement,
only the order-two case remains, except for the possible
order-seven or order-fourteen spectra over $K=\mathbb Q(\sqrt{-7})$.
An order-seven cubic spectrum is
$\{\zeta_7,\zeta_7^2,\zeta_7^4\}$ or its conjugate.
Either $x$ or $x^{-1}$ has age $(1+2+4)/7=1$.
In order fourteen the square has one of those order-seven
spectra, so an actual power is again a noncentral junior element.
Paragraph4 excludes both orders.  Thus all actual rank-three
elements over the seven fields are involutions.

For $K=\mathbb Q(i)$ the possible orders are two and four.
In order four the three roots are $i$ or $-i$.  Three equal
roots give age $3/4$ for $x$ or $x^{-1}$ and are excluded by
paragraph4.  The mixed spectra in (G31.2) remain, with minimum
age $5/4$ between the two inverse elements.  In every field
the involution has age $3/2$.  This proves the stated spectral
restriction without asserting that any of its surviving
possibilities occurs in a period-dominating cover.

\paragraph{6. Published low-degree route and the retained gate.}
For $15\le m\le35$ the rank-three element alone does not meet
$3\le\sqrt m/2$.  The precise published finite-list route is
Hiss--Malle, \emph{Low-dimensional representations of
quasi-simple groups}, LMS J.~Comput.~Math.~\textbf{4} (2001),
22--63, Theorem~1 on p.~22 and the conventions in Section~6, p.~30,
together with its complete corrected Table~2 in their
\emph{Corrigenda}, volume~\textbf{5} (2002), 95--126.
The correction includes previously omitted representations of
$L_3(3)$ in degrees16,26,27 and corrects the stated indicator
for a degree132 representation.  The intentionally omitted
alternating and $L_2(q)$ families must also be restored from
the original table conventions.  These are actual published
quasi-simple representation tables; using them for this gate
still requires the primitive-group structural reduction,
the action of every relevant outer normalizer, and the actual
$K$-model and inertia elements.  A bare list of complex
characters does not supply those maps.

The bounded-deviation Theorem~1.1 of the cited Guralnick--Tiep
paper is also insufficient by itself here.  For a rank-three
unitary element, $\dim V-|\operatorname{tr}(x)|\le6$,
so its $L^2$ deviation squared is at most12; the theorem's
bound $\max\{4C^2/63,40C\}$ at $C=12$ is480.
It does not eliminate dimensions15 through35.  The junior
Theorem~1.4 has been applied only through the actual
age-at-most-one elements in paragraphs4--5; it does not apply
to the surviving age-$3/2$ involutions or mixed Gaussian
age-$5/4$ elements merely from their moving rank.

Thus the remaining primitive branch has the finite range
(G31.1), retains tensor indecomposability, and has the exact spectra
above and the original period, full-group, field, tuple,
primitive-factor and lattice requirements retained.  HD-R12
further gives the field-dependent lower bounds stated at the
start.  The analysis does not assert that the absence of a
rank-one or rank-two element implies an absolute-irreducibility
or primitivity property for $N_3$; those are precisely the
branch addressed here.  The other rank-three branches are
handled separately in HD-S44 and subsequent work.

\paragraph{7. Removing the tensor hypothesis by the actual generator maps.}
HD-R14 proves directly that an absolutely irreducible finite group
in degree at least15 generated by actual moving-rank-three
elements preserves no nontrivial tensor decomposition, including
permutations of tensor factors.  Apply it to the original
$N_3\subset GL_K(V)$ after literal complexification.  The lower
bound $m\ge15$ above uses only the sphere-cover branch budget,
so this application does not depend on a tensor hypothesis.
Its contradiction is the explicit map
$w\mapsto\Phi^{-1}(w\otimes u)$ into a proper invariant
subspace after the original tensor-cycle and eigenvalue maps.
Consequently all conclusions of HD-G31 hold for the absolutely
irreducible complex-primitive $N_3$ branch without an additional
tensor assumption.  No new coefficient representation or
Jacobian factor is substituted in this application.

The supplementary exact program checks five noncoordinate
augmentation models, the scalar-cancelling commutator, the
finite $K$-intertwiner equations in two small models, primitive
matrix-unit adjoint identities, and the cancellation of the two
Schur scales.  It also factors every cyclotomic polynomial with
$\varphi(\ell)\le6$ over all eight fields, checks the actual
power ages, and checks the residual field intervals.  These are
finite exact checks of the displayed maps and arithmetic.  They
do not replace the published classifications or claim a finite
enumeration of the remaining primitive groups.
